\answer
\begin{enumerate}

%%--- 章末問題 3.1 ---
\item[\ref{ex3.1}]
$\exp(ev_D/kT) \gg 1$のとき$i_D \approx I_0\exp(ev_D/kT)$。
電流が10倍になる電圧増分$\mathit{\Delta}v$は
\begin{equation*}
\exp\left(\frac{e\mathit{\Delta}v}{kT}\right) = 10
\quad\Rightarrow\quad
\mathit{\Delta}v = \frac{kT}{e}\ln 10 = 0.026\times2.30 \approx 0.060\,\mathrm{V}
\end{equation*}
すなわちわずか60\,mVで電流は1けた変わる。
順方向電圧が0.6〜0.7\,V付近でほぼ一定に見えるのは，
この鋭い指数特性のためである。

%%--- 章末問題 3.2 ---
\item[\ref{ex3.2}]
順方向では$\exp(0.5/0.026) = \exp(19.2) \approx 2.2\times10^{8}$より$i_D \approx 2.2\times10^{8} I_0$，
逆方向では$\exp(-19.2) \approx 0$より$i_D \approx -I_0$。
比はおよそ2億倍で，たった1\,Vの向きの違いで電流が8けたも変わる。

%%--- 章末問題 3.3 ---
\item[\ref{ex3.3}]
MOSFETは式\ref{eq3.1}，IGBTは$P_{\mathrm{on}} = V_{\mathrm{on}}I$で見積もる。
$I = 10$\,AではMOSFET：$0.05\times10^2 = 5$\,W，IGBT：$1.25\times10 = 12.5$\,W。
$I = 50$\,AではMOSFET：$0.05\times50^2 = 125$\,W，IGBT：$1.25\times50 = 62.5$\,W。
損失は$R_{\mathrm{on}}I = V_{\mathrm{on}}$となる電流$I^* = V_{\mathrm{on}}/R_{\mathrm{on}} = 25$\,Aを境に大小が入れ替わる。
$I^2$で増える抵抗型は小電流に強く，$I$の1乗で増える電圧型は大電流に強い ―
「損失の質」の違いが素子の使い分けを生む（\ref{sec3.9}節）。

%%--- 章末問題 3.4 ---
\item[\ref{ex3.4}]
IGBTのゲートは酸化膜の上にあり，電場でp形表面に反転層を作って
「電子の供給路」を開け閉めする。
ゲート電圧を切れば反転層が消えて電子の供給が絶たれ，
それに伴い裏面からの正孔注入も止まるので，素子は自らオフできる
（内部の寄生サイリスタがラッチしない設計になっている）。

一方サイリスタのゲートは内側のp層に直接つながっており，
ゲート電流で2つのトランジスタの正帰還ループに火をつける。
いったんループが自立すると，
主電流のキャリア注入そのものが駆動電流を賄うため，
ゲートからはもはや介入できない。
「電場で通り道を作る」IGBTと「電流で正帰還を点火する」サイリスタ ―
同じ4層でも，制御のしくみの違いがオフの可否を分けている。

%%--- 章末問題 3.5 ---
\item[\ref{ex3.5}]
(a) 導通損失は
$P_{\mathrm{cond}} = V_{\mathrm{CE(sat)}}ID = 1.5\times100\times0.5 = 75$\,W。
スイッチング損失は式\ref{eq3.8}より
\begin{equation*}
P_{\mathrm{sw}} = \frac{1}{2}\times400\times100\times1\times10^{-6}\times10^4 = 200\,\mathrm{W}
\end{equation*}
総損失は$P_{\mathrm{total}} = 75 + 200 = 275$\,W。

(b) $P_{\mathrm{sw}}$は周波数に比例するから$200\times5 = 1000$\,Wとなり，
$P_{\mathrm{total}} = 1075$\,W。
損失は4倍近くに跳ね上がり，冷却が現実的でなくなる。
IGBTのインバータが数十kHz以下で運転される理由である。

(c) 導通損失は
$P_{\mathrm{cond}} = R_{\mathrm{on}}I^2D = 0.01\times100^2\times0.5 = 50$\,W。
スイッチング損失は
\begin{equation*}
P_{\mathrm{sw}} = \frac{1}{2}\times400\times100\times100\times10^{-9}\times5\times10^4
= 100\,\mathrm{W}
\end{equation*}
総損失は$150$\,Wで，(b)の約7分の1である。
切り替え時間が10分の1になったことで，
周波数を5倍に上げてなお損失を減らせる。
これがワイドギャップ半導体による高周波化の定量的な意味である。

%%--- 章末問題 3.6 ---
\item[\ref{ex3.6}]
\ref{sec3.9}節の手順（1.耐圧 → 2.電流 → 3.損失 → 4.駆動 → 5.熱）の順にたどる。

(a) \textbf{MOSFET（GaN HEMTを含む）}。
1.~耐圧：整流後の直流は約$141$\,V（入力変動を見て$200$\,V弱）だが，
ACアダプタは6章の絶縁型回路（フライバック）で，スイッチにはトランスの跳ね返り電圧が上乗せされる。
$600$〜$650$\,V級の定格なら50〜80\,\%の範囲に収まる。
2.~電流：$65$\,Wを$141$\,Vで割った平均電流は$0.5$\,A程度，ピークでも数Aで，
$I^* = V_{\mathrm{CE(sat)}}/R_{\mathrm{on}}$（十A級）よりはるかに小さい ― 抵抗性のMOSFETの領域である（\ref{fig3.9}）。
3.~損失：数百kHzでは切替え時間$1\,\mathrm{\mu s}$級のIGBTのスイッチング損失（式\ref{eq3.8}）が出力の半分に達し，使えない。
切替え時間数十nsのMOSFETなら数W，GaN HEMTならさらに小さい。
4.~駆動：電圧駆動で，ゲート電荷も小さく（GaNでは数nC），駆動回路は小さく損失も無視できる。
5.~熱：損失が$2$\,W程度なら，放熱器なしで基板の銅箔から逃がす熱抵抗（$30$\,K/W程度）でも
式\ref{eq3.11}の温度上昇は$60$\,Kで，筐体内が$50\,^{\circ}$Cでも$T_j \approx 110\,^{\circ}$Cと定格内に収まる。
小型化を重視するならGaN HEMTが最有力である。

(b) \textbf{SiC MOSFETまたはIGBT}。
1.~耐圧：電池$800$\,Vにサージを見込み，$1200$\,V級（$800/1200 \approx 67\,\%$）を選ぶ。
2.~電流：$150$\,kWを$800$\,Vで割ると$200$\,A近く，相電流は数百Aに達し，$I^*$を大きく超える ―
バイポーラ系（IGBT）の領域で，MOSFETを使うならチップを並列にする。
3.~損失：IGBTは$10$〜$20$\,kHzが限度で，章末問題\ref{ex3.5}の条件では1素子$275$\,W。
SiC MOSFETなら同じ耐圧でテール電流がなく，$50$\,kHzでも$150$\,Wと少ない。
4.~駆動：どちらも電圧駆動で，ゲート電荷が数百nCあっても駆動損失は$1$\,W以下にすぎない。
5.~熱：損失が百W級なので自然空冷では捨てられず，\ref{sec3.9}節の例のように
水冷で熱抵抗を$0.5$\,K/W程度まで下げて$T_j$を定格内に収める。
損失の少ないSiC MOSFETは，冷却器と電池を小さくできる分だけ車両全体で有利になる。

(c) \textbf{サイリスタ}。
1.~耐圧：数百kVは1素子の耐圧（$10$\,kV級が最大）をはるかに超えるので，
多数を直列にしたバルブとして使う。直列にできる素子であることが第一条件になる。
2.~電流：GW級を数百kVで割ってもkA級で，$I^*$を桁で超える ― 最大の電流容量をもつサイリスタの領域である。
3.~損失：商用周波数（$50$/$60$\,Hz）ではスイッチング損失は無視でき，損失はほぼ導通損失だけになる。
4.~駆動：電流駆動で，ゲートではオフできないが，商用周波数の交流に同期して電流が自然にゼロになるので他励形で運用できる。
高電位の素子へは光でゲート信号を送る。
5.~熱：1素子あたりの導通損失はkW級に達するので，純水で冷却して熱抵抗を数十分の1\,K/Wまで下げる。
電流駆動と転流の制約は，この電力規模では受け入れるしかない
（近年は，IGBTを多段に重ねてゲートでオフできるようにした自励式の変換所も増えている）。

\end{enumerate}
